% IEEE conference format (6-8 pages). Compile on Overleaf or with pdflatex.
% Figures are taken from ../results/figures (copy them next to this file on Overleaf).
\documentclass[conference]{IEEEtran}
\usepackage{graphicx}
\usepackage{booktabs}
\usepackage{amsmath}
\usepackage{hyperref}
\graphicspath{{../results/figures/}{figures/}}

\begin{document}

\title{From Scratch to Transfer Learning: An Empirical Study of CNN Design Choices on CIFAR-10}

\author{
\IEEEauthorblockN{Name 1, Name 2, Name 3}
\IEEEauthorblockA{Master 2 AI -- Programming for AI\\
Group X}
}

\maketitle

\begin{abstract}
% 150-200 words: problem, method, main numbers (baseline acc, best scratch acc, best transfer acc), main conclusion.
TODO
\end{abstract}

\begin{IEEEkeywords}
convolutional neural networks, data augmentation, transfer learning, fine-tuning, CIFAR-10
\end{IEEEkeywords}

\section{Introduction}
% Context, objectives (the 3 tasks), contributions, structure of the paper.
TODO

\section{Dataset and Pipeline}
% CIFAR-10, 45k/5k/10k split, normalization, DataLoaders, seed / reproducibility.
TODO

\section{Methodology}
\subsection{Baseline CNN}
% Architecture table (layer, output shape), number of parameters, receptive field,
% cross-entropy loss, Adam lr=1e-3, 20 epochs, best checkpoint on val accuracy.
TODO

\subsection{Controlled Experiments}
% One-factor-at-a-time protocol, list of variations (depth, kernel, stride/padding, BN,
% optimizers, learning rates, schedulers, augmentation, MixUp formula).
TODO

\subsection{Transfer Learning}
% ResNet-18 (ImageNet), frozen / finetune_layer4 / finetune_all, reduced lr, 128x128 inputs,
% ImageNet normalization, 10% vs 100% training data.
TODO

\section{Experiments and Results}
\subsection{Baseline}
\begin{figure}[t]
\centering
\includegraphics[width=\linewidth]{task1_curves.png}
\caption{Baseline CNN: training and validation loss and accuracy.}
\label{fig:baseline_curves}
\end{figure}

\begin{figure}[t]
\centering
\includegraphics[width=0.85\linewidth]{task1_confusion_matrix.png}
\caption{Baseline CNN: normalized confusion matrix on the test set.}
\label{fig:baseline_cm}
\end{figure}
% Discuss test accuracy, overfitting epoch, most confused classes.
TODO

\subsection{Architectural Variations and Hyperparameters}
\begin{table}[t]
\centering
\caption{Comparative benchmark (Task 2). Values from results/task2\_benchmark.csv.}
\label{tab:benchmark}
\begin{tabular}{lrrrrr}
\toprule
Run & Params & RF & Latency (ms) & Peak val acc & Gap \\
\midrule
Baseline & & & & & \\
Conv 2 / Conv 5 & & & & & \\
Kernel 5 / 7 & & & & & \\
Stride 2 / No padding & & & & & \\
BatchNorm & & & & & \\
SGD / AdamW & & & & & \\
lr 1e-2 / 1e-4 & & & & & \\
StepLR / Cosine & & & & & \\
Aug basic / full / MixUp & & & & & \\
\midrule
Best combined & & & & & \\
\bottomrule
\end{tabular}
\end{table}
% One paragraph per group with numbers + receptive field discussion.
TODO

\subsection{Transfer Learning}
\begin{table}[t]
\centering
\caption{Scratch vs frozen vs fine-tuned (Task 3).}
\label{tab:transfer}
\begin{tabular}{lrrrrr}
\toprule
Setup & Trainable & Ep. to 80\% & Test acc (10\%) & Test acc (100\%) & s/epoch \\
\midrule
Scratch & & & & & \\
Frozen & & & & & \\
Fine-tune layer4 & & & & & \\
Fine-tune all & & & & & \\
\bottomrule
\end{tabular}
\end{table}

\begin{figure}[t]
\centering
\includegraphics[width=\linewidth]{task3_convergence_100pct.png}
\caption{Validation curves of the four setups (100\% of the training data).}
\label{fig:convergence}
\end{figure}

\begin{figure}[t]
\centering
\includegraphics[width=0.85\linewidth]{task3_data_efficiency.png}
\caption{Data efficiency: test accuracy with 10\% and 100\% of the training data.}
\label{fig:data_efficiency}
\end{figure}
% Discuss data efficiency, convergence speed, final metrics, cost (params, time).
TODO

\section{Discussion}
% What worked and why (link to theory), limitations (single seed, 20 epochs, 128px), future work.
TODO

\section{Conclusion}
TODO

\begin{thebibliography}{1}
\bibitem{krizhevsky2009} A. Krizhevsky, ``Learning multiple layers of features from tiny images,'' Tech. Rep., 2009.
\bibitem{he2016} K. He, X. Zhang, S. Ren, and J. Sun, ``Deep residual learning for image recognition,'' in \emph{CVPR}, 2016.
\bibitem{ioffe2015} S. Ioffe and C. Szegedy, ``Batch normalization: Accelerating deep network training by reducing internal covariate shift,'' in \emph{ICML}, 2015.
\bibitem{kingma2015} D. P. Kingma and J. Ba, ``Adam: A method for stochastic optimization,'' in \emph{ICLR}, 2015.
\bibitem{loshchilov2019} I. Loshchilov and F. Hutter, ``Decoupled weight decay regularization,'' in \emph{ICLR}, 2019.
\bibitem{zhang2018} H. Zhang, M. Cisse, Y. N. Dauphin, and D. Lopez-Paz, ``mixup: Beyond empirical risk minimization,'' in \emph{ICLR}, 2018.
\bibitem{yosinski2014} J. Yosinski, J. Clune, Y. Bengio, and H. Lipson, ``How transferable are features in deep neural networks?'' in \emph{NeurIPS}, 2014.
\end{thebibliography}

\end{document}
